\problem{}

At noon-peak hour, a food-delivery platform still has $n$ unassigned orders.  
Each order can be delivered by exactly one of two on-line riders:

\begin{itemize}
\item Rider A (a new employee) handling order $i$ brings the platform a profit of $\text{profit1}[i]$ yuan;
\item Rider B (a veteran) handling order $i$ brings the platform a profit of $\text{profit2}[i]$ yuan.
\end{itemize}

To satisfy a company regulation, \textbf{Rider A must complete exactly $k$ orders} today, while the remaining orders are all assigned to Rider B.

If a total of $n$ orders are dispatched in this way, the platform's \emph{total revenue} is defined as the sum of the profits generated by the two riders:
\[
\text{Total Revenue} = \sum_{\text{orders to A}} \text{profit}_1[\cdot] + \sum_{\text{orders to B}} \text{profit}_2[\cdot].
\]

Return the maximum \emph{total revenue} the platform can obtain under the above requirement.

\noindent \textbf{Solution}:

\noindent \textbf{Step1}.
Create an array A\_better which A[i] = profit1[i] - profit2[i]. A\_better[i] represents the portion by which Rider A's profit exceeds Rider B's profit for order i.

\noindent \textbf{Step2}.
Sort A\_better in an ascending order.

\noindent \textbf{Step3}.
Calculate total revenue is the sum of the first k elements of A\_better plus the sum of all elements in array profit2 and return total return it. In fact, we assign the first k orders in A\_beetter to  rider A and the other orders to rider B.

\newpage